
\section*{PUERTA--LEY--01: taxonomía de puertas y rebote especular}

La corrección de nivel es: no hay que buscar una ``siguiente cifra'' sino una gramática de puertas. Cada puerta combina cuatro datos:
\[
\boxed{\text{capacidad }K(t),\quad\text{fase }K(t)\bmod9,\quad\text{poda/stutter},\quad\text{firma de frontera}.}
\]

La capacidad decimal viene dada por:
\[
\boxed{K(t)=\left\lfloor 6(t+1)\log_{1000}3\right\rfloor.}
\]
Cada bloque abre \(3^6=729\) subceldas; cada nueva triada poda por \(1000\). El calendario de poda no es periódico ordinario: es una palabra sturmiana de pendiente \(6\log_{1000}3\). Los stutters auditados son:
\[
21, 43, 65, 87, 109.
\]

\subsection*{1. Puertas iniciales}
\[
\begin{array}{c|c|c|c|c|c}
 t & K & K\bmod9 & (|\pi|,|e|,|\varphi|) & \#\text{orient} & \text{tipo}\\
\hline
0 & 0 & 0 & 729,729,729 & 5 & seed/sync/orientation \
1 & 1 & 1 & 378,422,349 & 3 & w30-lift/orientation \
2 & 2 & 2 & 238,368,388 & 2 & w30-lift/orientation \
3 & 3 & 3 & 184,284,173 & 5 & w30-lift/orientation \
4 & 4 & 4 & 207,207,122 & 4 & w30-lift/orientation \
5 & 5 & 5 & 39,151,151 & 1 & R36-EF/unique-local \
6 & 6 & 6 & 110,110,69 & 1 & unique-local \
7 & 7 & 7 & 81,29,81 & 3 & orientation \
8 & 8 & 8 & 59,59,59 & 3 & sync/orientation \
9 & 9 & 0 & 43,43,43 & 2 & sync/orientation \
10 & 10 & 1 & 32,32,32 & 1 & sync/unique-local \
11 & 11 & 2 & 23,24,23 & 1 & unique-local \
12 & 12 & 3 & 17,17,17 & 1 & sync/unique-local \
13 & 13 & 4 & 13,13,13 & 1 & sync/unique-local \
14 & 14 & 5 & 9,10,9 & 1 & unique-local \
15 & 15 & 6 & 7,7,8 & 1 & unique-local \
16 & 16 & 7 & 6,6,6 & 1 & sync/unique-local \
17 & 17 & 8 & 4,4,4 & 1 & sync/unique-local \
18 & 18 & 0 & 4,4,4 & 1 & sync/unique-local \
19 & 19 & 1 & 3,3,3 & 1 & sync/unique-local \
20 & 20 & 2 & 2,2,2 & 1 & sync/unique-local \
21 & 20 & 2 & 611,527,723 & 5 & stutter/orientation \
22 & 21 & 3 & 421,242,483 & 2 & orientation \
\end{array}
\]

La ventana inicial muestra tres hechos: \(w_6\to w_{30}\) ocupa \(t=0,\ldots,4\); la frontera no trivial \(R36\) empieza en \(t=5\); y el paquete \(t=7,8,9\) es una puerta de orientación.

\subsection*{2. Puertas de orientación}
En una puerta de orientación la firma deja varias matrices locales, pero todas tienen:
\[
\boxed{\varphi\text{ fijo},\qquad \pi+e\text{ constante}.}
\]
La ambigüedad vive sólo en el reparto espejo \(\pi/e\), y la supervivencia futura selecciona una orientación.

\[
\begin{array}{c|c|c|c|c|c|c|c}
 t & K & K\bmod9 & (|\pi|,|e|,|\varphi|) & \# & \varphi & \pi+e & \text{surv real}\\
\hline
0 & 0 & 0 & 729,729,729 & 5 & 121200 & 211012 & 3442951 \
1 & 1 & 1 & 378,422,349 & 3 & 112202 & 100110 & 1331000 \
2 & 2 & 2 & 238,368,388 & 2 & 121020 & 112222 & 531441 \
3 & 3 & 3 & 184,284,173 & 5 & 010210 & 011000 & 205379 \
4 & 4 & 4 & 207,207,122 & 4 & 010200 & 212202 & 79507 \
7 & 7 & 7 & 81,29,81 & 3 & 200221 & 101001 & 4913 \
8 & 8 & 8 & 59,59,59 & 3 & 110110 & 112000 & 2197 \
9 & 9 & 0 & 43,43,43 & 2 & 010122 & 120212 & 810 \
21 & 20 & 2 & 611,527,723 & 5 & 211021 & 211122 & 7723188 \
22 & 21 & 3 & 421,242,483 & 2 & 101021 & 202011 & 3027600 \
24 & 23 & 5 & 371,371,332 & 2 & 210002 & 222220 & 456533 \
26 & 25 & 7 & 184,197,198 & 6 & 121011 & 112212 & 72324 \
28 & 27 & 0 & 106,106,105 & 5 & 021112 & 211002 & 11638 \
30 & 29 & 2 & 56,56,57 & 2 & 211110 & 020011 & 1728 \
31 & 30 & 3 & 18,41,38 & 2 & 000011 & 202022 & 1000 \
43 & 41 & 5 & 425,586,697 & 3 & 111102 & 002101 & 6751269 \
44 & 42 & 6 & 654,459,408 & 3 & 102200 & 012111 & 2628072 \
45 & 43 & 7 & 400,486,226 & 3 & 211122 & 002100 & 1030301 \
46 & 44 & 8 & 355,354,354 & 5 & 102022 & 102020 & 405224 \
47 & 45 & 0 & 258,113,246 & 6 & 022110 & 121210 & 157464 \
48 & 46 & 1 & 146,189,189 & 2 & 020122 & 110112 & 62400 \
50 & 48 & 3 & 101,81,101 & 5 & 111221 & 100112 & 10648 \
52 & 50 & 5 & 54,54,54 & 3 & 011021 & 021121 & 1728 \
53 & 51 & 6 & 39,40,40 & 3 & 020002 & 210211 & 576 \
65 & 62 & 8 & 546,698,729 & 3 & 012222 & 022002 & 5864400 \
66 & 63 & 0 & 197,527,141 & 3 & 010112 & 100022 & 2299968 \
67 & 64 & 1 & 119,399,452 & 10 & 011120 & 011202 & 884736 \
69 & 66 & 3 & 247,197,247 & 5 & 200222 & 020111 & 140608 \
70 & 67 & 4 & 180,180,175 & 2 & 221000 & 202100 & 52022 \
72 & 69 & 6 & 96,96,96 & 6 & 122210 & 112002 & 8000 \
\end{array}
\]

\subsection*{3. La puerta nueve}
Hay dos lecturas coherentes de la novena puerta.

Noveno bloque, \(t=8\):
\[
|\pi|=|e|=|\varphi|=59,
\]
la firma deja tres orientaciones con \(\varphi=110110\) fijo y \(\pi+e=112000\). La supervivencia elige la rama central.

Novena triada decimal, \(t=9\):
\[
|\pi|=|e|=|\varphi|=43,
\]
la firma deja dos orientaciones con \(\varphi=010122\) fijo y \(\pi+e=120212\). La supervivencia elige la primera.

\[
\boxed{\text{la novena puerta es una puerta de orientación: }\varphi\text{ fija escala y }\pi/e\text{ se espejan.}}
\]

\subsection*{4. Por qué el diez no es una repetición literal del uno}
La puerta 10 no repite literalmente la puerta 1 en conteos ni firmas. Lo que se reinicia es la fase digital:
\[
K=9\Rightarrow K\bmod9=0,
\qquad
K=10\Rightarrow K\bmod9=1.
\]
Como \(1000\equiv1\pmod9\), el retorno \(10\to1\) es retorno de fase \(dr_9\), no igualdad de frontera. La frontera queda transformada por la memoria de carry y supervivencia.

\subsection*{5. Estadística por fase}
\[
\begin{array}{c|c|c|c|c}
K\bmod9 & \#\text{capas} & \#\text{sync} & \#\text{orient} & \#\text{stutter}\\
\hline
0 & 12 & 5 & 5 & 0 \
1 & 12 & 6 & 3 & 0 \
2 & 14 & 4 & 5 & 2 \
3 & 12 & 3 & 6 & 0 \
4 & 12 & 4 & 2 & 0 \
5 & 15 & 5 & 4 & 2 \
6 & 13 & 2 & 4 & 0 \
7 & 13 & 5 & 4 & 0 \
8 & 13 & 6 & 6 & 1 \
\end{array}
\]

\subsection*{6. Ley de puerta}
La taxonomía ya no es plana. Una puerta puede ser:
\[
\boxed{\text{poda}},\quad
\boxed{\text{stutter/reapertura}},\quad
\boxed{\text{orientación especular}},\quad
\boxed{\text{saturación}},\quad
\boxed{\text{neutralidad dual}}.
\]
La ley general debe leerse como un autómata de frontera:
\[
\boxed{\text{abrir }729\to\text{poda/cilindro}\to\text{firma}\to\text{orientación}\to\text{supervivencia}.}
\]

\subsection*{Dictamen}
\[
\boxed{\text{verde: la puerta 9 pertenece a una familia de orientación, no es anécdota.}}
\]
\[
\boxed{\text{verde: los stutters forman el calendario de reapertura por la ley }729/1000.}
\]
\[
\boxed{\text{verde: el retorno 10=1 es retorno de fase }K\bmod9\text{, no repetición literal de la frontera.}}
\]
\[
\boxed{\text{amarillo: falta convertir esta taxonomía en un transductor EF automático completo.}}
\]
